\documentclass[11pt]{article}
\usepackage[margin=1in]{geometry}
\usepackage[T1]{fontenc}
\usepackage[utf8]{inputenc}
\usepackage{lmodern}
\usepackage{booktabs}
\usepackage{longtable}
\usepackage{array}
\usepackage{graphicx}
\usepackage{fancyhdr}
\usepackage{hyperref}
\usepackage{xcolor}

\graphicspath{{../../}}

\hypersetup{
  colorlinks=true,
  linkcolor=blue,
  urlcolor=blue,
  pdftitle={From ERS-111 Behavior Diagrams to Distributed Embodiment},
  pdfauthor={Dr. Christopher A. Tucker}
}

\pagestyle{fancy}
\fancyhf{}
\fancyfoot[C]{\thepage}
\renewcommand{\headrulewidth}{0pt}
\renewcommand{\footrulewidth}{0pt}

\title{From ERS-111 Behavior Diagrams to Distributed Embodiment:\\SOFTSIM and GA144 as an Executable Framework}
\author{Dr.\ Christopher A.\ Tucker}
\date{}

\begin{document}
\maketitle

\begin{abstract}
This paper proposes a follow-on implementation framework derived from
prior analysis of Sony \texttt{ERS-111} \texttt{R-CODE} behavior
diagrams. The earlier corpus study showed that many sample routines are
organized around a compact embodied state vocabulary centered on
initialization, sensing, iterative action, synchronization, and
recovery. The present paper asks how such behavioral organization may
be carried toward executable distributed control. Its central claim is
that behavior diagrams can function not only as descriptive artifacts,
but also as executable organizational models when they are mapped onto
communicating node ensembles, initialized through the documented
GreenArrays workflow, and exercised functionally in \texttt{SOFTSIM}.
In this framing, \texttt{SOFTSIM} serves as a functional execution and
inspection environment, while \texttt{GA144} provides a candidate
distributed asynchronous hardware target for studying fast reactive
loops, slower supervisory intervention, and low-energy embodied control
organization.
\end{abstract}

\section*{Introduction}

This paper is a follow-on to the \texttt{ERS-111} behavior-diagram
study rather than a replacement for it. The earlier paper was closed
around the corpus itself: it aggregated the generated
\texttt{R-CODE} diagrams, identified their recurring state vocabulary,
and argued that apparently different routines are often built from a
compact embodied grammar. That result remains the starting point here,
but not the final objective.

The present question is implementation-oriented. If a robot behavior
corpus can be described as a recurring state vocabulary, can that same
vocabulary support a path toward executable distributed control? More
specifically, can behavior diagrams be treated as organizational models
for implementation on a communicating many-node substrate rather than
only as retrospective diagrams of legacy code?

This question matters because current robotics research increasingly
distinguishes between faster embodied control and slower supervisory or
reasoning layers. In contemporary work, low-latency action generation
and higher-level semantic or deliberative guidance are often separated
functionally even when they remain tightly coupled in deployment
\textit{(Zhu et al., 2024; Bjorck et al., 2025; Song et al., 2025)}.
The present work does not claim priority over that general idea.
Instead, it explores a different contribution: how a compact
behavioral vocabulary may be instantiated, inspected, and eventually
realized on a distributed asynchronous substrate.

The paper makes four main contributions. First, it restates the
\texttt{ERS-111} corpus result in the narrower form needed for
implementation work: the diagrams expose a compact embodied state
vocabulary rather than a collection of unrelated scripts. Second, it
shows that the GreenArrays documentation provides a continuous
toolchain from per-node implementation through declarative
initialization to functional simulation. Third, it argues that
behavior diagrams can be treated as executable organizational models
for distributed control rather than only as analytic artifacts.
Fourth, it defines a conservative experiment program in which
\texttt{SOFTSIM} is used for functional execution and inspection while
\texttt{GA144} remains the hardware target for later realization rather
than a source of premature timing claims.

\section*{Background from ERS-111}

The prior \texttt{ERS-111} study matters here because it supplies the
abstraction that this paper seeks to operationalize. At corpus scale,
the generated \texttt{R-CODE} diagrams repeatedly returned to a small
set of state-types such as \texttt{Sense / Decide},
\texttt{Action Loop}, \texttt{Boot / Safe Pose},
\texttt{Synchronize}, \texttt{Sense Fall State}, and
\texttt{Recover}. The significance of that finding was not merely that
one title appeared more frequently than another. The stronger result
was that many distinct routines were organized through repeated
recombination of a compact reactive and embodied control vocabulary.

That result provides a useful bridge for implementation work. If
behavior can be described at the level of recurring compartments such
as initialization, sensing, bounded action, synchronization, and
recovery, then those compartments can in principle be mapped onto
concrete execution units. The present paper therefore uses the
\texttt{ERS-111} corpus not as its primary object of analysis, but as
the source of a compact behavioral grammar that may be carried forward
into implementation experiments.

\section*{GreenArrays Toolchain}

The GreenArrays documentation supplies the implementation path required
for this next step. \texttt{DB001} describes the \texttt{F18A} node as
a highly constrained asynchronous processing element with local memory,
direct port execution, and communication semantics centered on
suspension and handoff across neighboring nodes. This matters because
it means that distributed behavior is not added as an afterthought.
Communication is part of the machine model itself.

\texttt{DB013 Chapter 4} shows how code is assembled per node and
stored through the object-bin workflow. That chapter is important
because it demonstrates that the programmer is not targeting a single
flattened image alone, but a persistent collection of node-local
artifacts. \texttt{DB013 Chapter 6} then shows how those artifacts can
be gathered into a Boot Descriptor Language workflow that defines how a
multi-node configuration is initialized and streamed into execution.

\texttt{DB013 Chapter 7} completes the path by documenting
\texttt{SOFTSIM} as a high-level software simulator for up to two
\texttt{GA144} chips. Its role is not to reproduce physical timing
faithfully, but to provide functional execution, operator inspection,
stepping, and breakpoint control for distributed node programs. That
distinction is critical. The present paper therefore uses
\texttt{SOFTSIM} as a functional and observational environment, not as
a timing-accurate proof of hardware behavior.

Taken together, these chapters define a continuous workflow:

\begin{center}
\texttt{node implementation -> object bins -> BDL initialization -> SOFTSIM execution -> GA144 realization}
\end{center}

That continuity is one of the main reasons the GreenArrays environment
is relevant to the present paper. It offers a concrete bridge from
behavioral organization to executable distributed embodiment.

\section*{Behavior Diagrams as Executable Models}

The central interpretive claim of this paper is that behavior diagrams
can be promoted from descriptive visualizations into executable
organizational models. A state diagram already distinguishes stable
behavioral compartments, entry conditions, transition conditions, and
return paths. These distinctions are analytically useful on their own,
but they are also structurally close to what distributed implementation
requires.

On a node-based asynchronous substrate, one can treat each state as a
small executable routine with bounded responsibilities. Some states may
remain local to a single node. Others may be split across
communicating nodes when sensing, local action, and supervisory
intervention should be kept distinct. In this view, the transition
structure functions as an implementation guide: it suggests where
behavior must remain local, where communication must occur, and where
synchronization or recovery paths need to be observable.

This is not yet a claim that every behavior diagram should be mapped
one-to-one onto hardware nodes. It is a more modest claim. The
diagrams provide a tractable intermediate representation that can guide
distributed implementation decisions while preserving behavioral
clarity. They help identify what is stable in the organization of a
behavior even when the sensing words, actuation words, or supervisory
policies later change.

\section*{Fast and Slow Cadence}

Recent robotics literature gives added relevance to this approach by
showing that explicit fast/slow or dual-system decompositions are now
common enough to support comparison. In \texttt{GR00T N1}, \texttt{Hume},
and \texttt{Language-Conditioned Robotic Manipulation with Fast and
Slow Thinking}, faster reactive or action-generation layers are paired
with slower supervisory, semantic, or reasoning-capable layers
\textit{(Zhu et al., 2024; Bjorck et al., 2025; Song et al., 2025)}.

The present work does not claim that the GreenArrays environment
already matches the scale or capability of those systems. Its value is
different. It offers a small, inspectable, distributed implementation
framework in which a similar cadence can be studied concretely. In the
working interpretation used here, fast reactive behavior corresponds to
local low-latency node execution, while slower supervisory behavior
corresponds to a less frequent layer that modulates or redirects local
loops.

Within that framing, \texttt{GA144} is valuable not because it replaces
contemporary large-scale robot foundation models, but because it offers
a contrasting substrate for embodied control experiments. As a
distributed asynchronous many-node device, it supports organizing
behavior as multiple small communicating reactive processes rather than
as one monolithic controller. This is especially relevant for the
present work because it provides a plausible low-energy hardware target
for the reactive layer while \texttt{SOFTSIM} provides the functional
environment needed to inspect node state, communication, and
supervisory handoff before stronger hardware claims are made.

\section*{Experiment Program}

The experiment program proposed here is deliberately conservative. Its
purpose is not to demonstrate superiority over current embodied AI
systems, nor to make timing claims that \texttt{SOFTSIM} cannot
support. Its purpose is to establish a credible sequence of executable
evidence.

The first paper-facing target in that sequence should be the existing
\texttt{Move} diagram from the \texttt{ERS-111} study. \texttt{Move}
is compact, reactive, and already organized around a loop of repeated
action, fall sensing, recovery, and return to execution. That makes it
the best first candidate for a direct diagram-to-program mapping using
the documented \texttt{1737}-style minimal load path in
\texttt{SOFTSIM}. The immediate implementation goal is not yet a full
locomotion controller, but a minimal structural mutation of the
\texttt{1737} baseline whose repeated region, alternate region, and
return-to-loop behavior correspond to the main organizational features
of \texttt{Move}. A second small target should then be a reduced
recovery-oriented control slice centered on \texttt{Boot / Safe Pose},
\texttt{Sense Fall State}, \texttt{Recover}, and \texttt{Action Loop},
rather than a more specialized family such as \texttt{Football}.

\begin{itemize}
\item establish that the local \texttt{SOFTSIM} installation is a usable functional execution environment
\item preserve a smallest confirmed baseline such as the documented \texttt{1737} example
\item demonstrate a minimal one-node mutation that proves custom executable behavior is possible
\item demonstrate a minimal two-node reactive/supervisory handoff
\item make node-local state and communication markers observable during execution
\item carry at least one stable organization toward later realization on \texttt{GA144}
\end{itemize}

This sequence matters because it distinguishes functional evidence from
hardware-performance evidence. If the sequence succeeds in
\texttt{SOFTSIM}, it supports claims about executable organization,
inspection, and distributed behavioral decomposition. If a later
hardware realization succeeds on \texttt{GA144}, it supports the
stronger claim that at least one such organization survives translation
onto the physical substrate. It still would not by itself prove
large-scale robot-level capability, but it would validate the
implementation pathway.

\section*{Claim Boundaries}

Because this work sits between historical analysis and future hardware
implementation, its claim boundaries must remain explicit.

\textbf{What can be claimed safely.} The prior \texttt{ERS-111} corpus
study identified a compact embodied state vocabulary. The GreenArrays
documentation defines a continuous path from per-node implementation
through declarative initialization to functional simulation. Live local
execution confirms that \texttt{SOFTSIM} can load, step, and inspect
distributed node programs. These points support the claim that behavior
diagrams can function as executable organizational models in a
distributed asynchronous implementation workflow.

\textbf{What should not yet be claimed.} The present evidence does not
show timing fidelity in \texttt{SOFTSIM}, does not demonstrate
superiority to contemporary foundation-model robot architectures, and
does not prove robot-level performance on real embodied tasks. It also
does not show that every behavior diagram admits a simple one-to-one
mapping onto node hardware.

These boundaries do not weaken the contribution. Rather, they define it
more clearly. The paper's value lies in identifying and operationalizing
a tractable bridge between behavioral abstraction and executable
distributed embodiment.

\section*{Conclusion}

The original \texttt{ERS-111} behavior-diagram paper established that
the Sony \texttt{R-CODE} corpus is organized around a compact embodied
state vocabulary. The present follow-on paper asks what that result can
be used for. Its answer is that such a vocabulary can serve as an
intermediate representation between historical behavior diagrams and
distributed executable control.

The GreenArrays environment is especially useful for this purpose
because it provides a continuous path from node-local implementation,
through declarative multi-node initialization, to functional execution
and inspection in \texttt{SOFTSIM}, and ultimately to realization on
\texttt{GA144}. In that sense, the contribution is not the invention of
fast/slow embodied control as an abstract concept. It is the proposal
of an executable framework for studying how such control may be
organized, inspected, and carried toward hardware on a distributed
asynchronous substrate.

\section*{References}

\begin{itemize}
\item Arkin, Ronald C. \textit{Behavior-Based Robotics}. MIT Press, 1998.
\item Bjorck, Johan, Eric Jang, Dibya Ghosh, et al. ``GR00T N1: An Open Foundation Model for Generalist Humanoid Robots.'' arXiv, 2025. \url{https://arxiv.org/abs/2503.14734}
\item Brooks, Rodney A. ``Intelligence without Representation.'' \textit{Artificial Intelligence} 47, no. 1-3 (1991): 139-159.
\item Song, Jiankai, Yilun Du, Kalyan S. Perumalla, et al. ``Hume: Toward a Universal and Fine-Grained Human-Centric Understanding Model for Embodied AI.'' arXiv, 2025. \url{https://arxiv.org/abs/2505.21432}
\item Xu, Zhiyuan, Kun Wu, Junjie Wen, Jinming Li, Ning Liu, Zhengping Che, and Jian Tang. ``A Survey on Robotics with Foundation Models: toward Embodied AI.'' arXiv, 2024. \url{https://arxiv.org/abs/2402.02385}
\item Zhu, Xianghao, Yizhou Wang, Junheng Hao, et al. ``Language-Conditioned Robotic Manipulation with Fast and Slow Thinking.'' arXiv, 2024. \url{https://arxiv.org/abs/2401.04181}
\end{itemize}

\end{document}
